\documentclass[11pt,a4paper]{article}
\usepackage[margin=26mm]{geometry}
\usepackage{amsmath,amssymb,amsthm,hyperref,longtable,booktabs,array}
\emergencystretch=1em
\newcommand{\Sz}{\operatorname{Sz}}
\title{Subfield descent and the higher-rank Suzuki constructions}
\author{Human Compression project}\date{}
\begin{document}\maketitle
\begin{abstract}
We account for the mathematical claims in the fixed twenty-page edition
of \emph{The residually weakly primitive pre-geometries of the Suzuki
simple groups}. Its exhaustive theorem is false. We retain the earlier
counterexample and give proofs or explicit corrections for its remaining
claims. The surviving construction families follow from subfield descent,
cyclic coset coordinates and scalar irreducibility. Several other prescribed
intersection lattices are impossible, and a scalar codimension test rules
out the mixed construction. This is a complete treatment of the fixed
source's assertions, not a corrected classification of all Suzuki geometries.
\end{abstract}

\section{The source and the meaning of completion}
The source is D.~Leemans, Note di Matematica 20(1) (2000/2001),
1--20, received 26 July 1999 and accepted 21 September 1999.
We use the retained published edition, whose physical and printed page
numbers agree. The exact source file and page references are in the
claim ledger. The companion volume includes all the proof texts used
below, and its acquisition record binds their earlier independent reviews.

For $G_i\le G$ and a finite type set $I$, write
$G_J=\bigcap_{j\in J}G_j$ and $G_\varnothing=G$.
The source's condition is
\[
 \text{for every }J\subsetneq I\text{ there is }i\notin J
 \text{ with }G_{J\cup\{i\}}\text{ maximal proper in }G_J.
 \tag{1}
\]
The source uses left cosets and left translation. Its residues for flags
inside the distinguished chamber are the algebraic intersection systems.
Its local two-transitivity condition applies to those algebraic rank-one
residues. Actual thinness counts all chambers incident with an actual panel.
The distinction is necessary for the corrections imported from the
rank-three treatment.
On p.~5 the source defines groups of Suzuki type to be simple and
explicitly excludes $\Sz(2)$. We apply that convention to the Suzuki
labels in its residue lists. In particular the common Suzuki bottom in
$R_{10},R_{11},R_{13}$ is simple; the parameter is at least eight.


We use flag transitivity in the source's stated strong sense, including
transitivity on every flag type. In the positive results here we prove
that every flag has a common representative in $G$, which implies both
flag transitivity and the geometry axiom. No inference from chamber
transitivity alone is needed.

The source explicitly states on p.~3 that its principal theorem concerns
pregeometries, without assuming flag transitivity or the geometry axiom.
On p.~15 it requires the \emph{whole intended intersection lattice} of
each construction, as inferred from the corresponding rank-four case.
A list of abstract parabolic isomorphism types is therefore weaker than
the construction being asserted. We keep these two meanings separate.
A refuted exhaustive theorem is disposed of by a counterexample; it does
not create an obligation to prove a stronger repaired classification.

\section{Common theory and the rank-three input}
The classical premises are the four-dimensional Suzuki model, its
maximal-subgroup and centralizer structure, simple subfield maximality,
Lang--Steinberg surjectivity and Schur--Zassenhaus. None is presented as
new. The source itself cites L\"uneburg for the subgroup structure.
The written common theory in the volume supplies the needed matrix
intersections, torus arithmetic, normalizer restrictions and affine
uniqueness arguments. The source's own exhaustive rank-three list is
never used as a premise.

Lemma~1, that every rank-three coset pregeometry is a geometry, has the
short common-representative proof in the acquired preliminary paper.
Lemma~2 is the odd-exponent gcd identity for $2^{2m}+1$.
Lemma~3 is the proper-overgroup statement for an embedded full split
dihedral in a maximal simple subfield; the ambient group itself must be
excluded from its literal word ``subgroups.'' Lemma~4 follows from the
proved exclusion of the other maximal subgroup types in rank at least
four. The consequences of (1) on pp.~4--5, including strict intersections
and a maximal chain obtained by ordering labels, are proved in the common
prime-degree argument.

Theorems~1 and 2 are recovered: the three necessary prime-degree patterns
are proved without an enumeration assumption, each allowed pattern has
an explicit realization, and a prime field degree with composite $q-1$
admits no rank at least three satisfying (1). The thirteen labels in
Theorem~3 have the following exact relation to the rank-three treatment.
\begin{center}
\begin{tabular}{lll}
\toprule
L110 labels & Rank-three source labels & Disposition\\
\midrule
$R_1,R_2,R_3$ & Theorem 5.3, patterns 1--3 & necessary and realizable\\
$R_4,R_5,R_6$ & Theorem 5.5, patterns 1--3 & prescribed meets impossible\\
$R_7$ & Theorem 5.5, pattern 4 & canonical realization\\
$R_8$ & Theorem 5.5, pattern 5 & prescribed meet impossible\\
$R_9,R_{10},R_{11}$ & Theorem 5.5, patterns 6--8 & canonical realizations\\
$R_{12}$ & Theorem 5.5, pattern 9 & bare meets exist; (1) fails\\
$R_{13}$ & Theorem 5.5, pattern 10 & canonical realization\\
\bottomrule
\end{tabular}
\end{center}
The exhaustive assertion in Theorem~3 is false: the acquired
$\Sz(8),\Sz(32),C_{533}:C_4$ family in $\Sz(32768)$ has
intersection orders $20,52,164,4$, whereas the applicable printed pattern
requires the two simple parabolics to meet in order eight. Its geometry,
firmness, residual connectedness and flag transitivity are proved.
The exact corrections, including the scalar obstruction to $R_{12}$,
are imported with their proofs rather than rederived by a new search.

Theorem~4 is explicitly a quotation of the separate dihedral classification,
reference [20]. That original is outside the fixed accessible intake.
We retain the quotation as an attributed premise, without claiming recovery
of the cited original. The needed odd-dihedral maximality consequences
also have elementary proofs in the common theory and the prime-residue
extension paper; the higher constructions do not rely on an unchecked
machine enumeration of dihedral geometries.

\section{The rank-four extension argument}
The new nonextension paper proves that none of $R_1,R_2,R_3$ extends
to a rank-four system with its simple parabolic maximal in the ambient
Suzuki group. For $R_1$, the affine subgroup forces its containing
soluble parabolic into a point stabilizer, which cannot contain the
required odd dihedral. For $R_2$, the common-reflection argument aligns
the two prime-degree simple subfields and makes their meet $\Sz(2)$,
of order twenty, not the prescribed dihedral of order ten.
For $R_3$, the common full torus forces an affine meet, and the unique
simple subfield containing that embedded affine group makes two
parabolics equal. Equal orders alone would not justify that last step.
Thus the source's claimed surviving $R_2$ candidate and Figure~1 are
refuted, while the nonextension conclusions for $R_1,R_3$ are recovered.
Section~3.2 follows from the recovered prime-degree obstruction.

In Section~3.3, the $R_4,R_5,R_6,R_8$ extension branches already start
from impossible full residue lattices. Their higher-rank continuations
cannot repair the obstruction. The $R_{12}$ branch fails (1) by the
scalar codimension test. The other five branches require a separate
necessity argument: existence of the canonical families does not prove
that every extension of a prescribed residue has their lattice.

The split extension paper proves this necessity for $R_7$ and $R_9$.
A field-generating cyclic complement identifies the embedded subfield
inside every simple overgroup of the prescribed dihedral or central
affine subgroup. Its intersection with the maximal parabolic then
excludes such a simple overgroup. The remaining dihedral parabolic lies
in the ambient split normalizer; the affine parabolic lies in the
appropriate point stabilizer. Scalar irreducibility and the absence of
intermediate overgroups force the full central radical and the same
cyclic complement in both soluble parabolics. The last cyclic and
dihedral maximality conditions force two distinct primes and every
remaining intersection. This proves the conditional extension statement
also in the repeated $R_7$ branch of Section~3.3.2. On p.~13 the stray
$G_4$ means $G_3$. The printed divisibility for the new prime must be
$n'\mid(s'-1)/(t-1)$; the quotient $(q-1)/(s'-1)$ is not the correct
condition. These textual corrections do not weaken the necessity claim.

For $R_{10}$ and $R_{11}$, use $H=C_m:C_c$ with the full torus
complement, of order two in the split case and four in the nonsplit
case. Its normalizer is $H$ itself. Conjugacy inside each subfield group
therefore makes the containing subfield of each degree unique. Two
such groups containing $H$ meet in the compatible gcd field; in the
nonsplit argument the full rotation order exceeds five. We do not
assume that the bottom torus of order five determines $\Sz(8)$.
A divisor argument supplies the remaining information: an intersection
system of $r$ cyclic extensions with a common complement has $r$
prime-index steps, and strictness supplies $r$ different prime valuation
witnesses. Its top rotation order divided by its bottom order is thus
a product of $r$ distinct primes, with one omitted at each parabolic.
This remains true when a selected prime divides the bottom order.

The torus necessity paper also derives the input field degrees from
(1). A composite relative degree would insert an intermediate compatible
torus order, contradicting the prime quotient unless the simple bottom
were maximal, which is impossible at that composite degree. Applying the
same argument in the ambient residue gives the two distinct relative
primes and their product. The source's $R_{11}$ row had not stated those
degree identities explicitly. With these premises proved, the fixed
maximal simple parabolic determines the other two simple parabolics,
and the full-complement and divisor arguments determine the remaining
torus parabolic and all its intersections. The $R_{10}$ paragraph's
later interchange of $G_2,G_3$ is corrected to agree with its opening
lattice. For $R_{13}$ the common simple bottom and uniqueness align all
field groups; maximal prime degrees then give every intersection.

These are conditional necessities for the exact source residues.
They do not make the source's global branch list exhaustive. The first
assembly conflated these two questions; its independent review held
completion on that gap. That assembly and hold are preserved. The two
new necessity papers supply the missing arguments in this edition,
while the global exhaustive conclusions remain refuted.

\section{Every numbered construction}
Put $S_d=\Sz(2^d)$. In the positive families choose distinct odd primes
$a_i$, an odd $e>1$, $t=2^e$, $n=e\prod_i a_i$, and standard
field parabolics $S_{n/a_i}$. Selecting a set $J$ of field labels leaves
$S_{n_J}$ with $n_J=e\prod_{i\notin J}a_i$.
The following table states the complete dispositions; the accompanying
papers give all subgroup and residue proofs.
\begin{longtable}{p{0.08\linewidth}p{0.84\linewidth}}
\toprule
No. & Disposition of the source recipe\\\midrule\endhead
1 & Its intended full lattice contains the forbidden shared-reflection
prime-subfield meet. It is impossible; the additional affine parabolic
cannot repair it.\\
2--4 & The intended lattices contain the impossible $R_4,R_5,R_6$
residues. Separately, their literal full intersections of orders 10, 4,
4 cannot be subgroups of the displayed odd dihedral: its rotation order
has no factor 5 and its two-part is 2. Deleting the last index would
repair that notation, but would not repair the intended lattice.\\
5 & After restoring the missing factor $t-1$ in the dihedral order and
correcting the conflicting indices and the undefined product index,
the canonical split family satisfies (1) for every allowed parameter.
The two cases have $e$ prime, including the repeated-prime case $a_i=e$.
If several distinguished primes give valid choices, different omitted
simple-subfield orders distinguish the resulting stated-action families.\\
6 & The intended $R_8$ lattice is impossible. This does not exclude the
bare parabolic types: the retained example with relative degrees
$3,5,7$ and torus primes $13,41,113$ realizes those types with a different
lower lattice, as a thick, residually connected, flag-transitive geometry.\\
7 & Relabel the repeated subfield/dihedral index. The common scalar field
lemma proves all intersections and (1) for every allowed parameter.\\
8--9 & Relabel the repeated final subfield/torus index. Every allowed
parameter has a canonical realization satisfying (1), thickness,
residual connectedness and flag transitivity. Selected primes may divide
the base torus order; the valuation proof includes this case.\\
10 & Relabeling the repeated final index does not repair the mathematics.
The bare intended intersections exist. For every embedding with the
prescribed mixed parabolic and pair intersections, the remaining
intersections are all nonmaximal in that parabolic, so (1) fails.\\
11 & Every literal allowed parameter has a canonical thick, residually
connected, flag-transitive realization satisfying (1). The exact
criterion is $e>1$, which also permits all distinct prime divisors to be
selected when a repeated prime factor remains. That boundary is absent
from the source rank bound.\\\bottomrule
\end{longtable}

For Constructions~5 and 7 put $M=(t-1)\prod_i\ell_i$, where the
$\ell_i$ are distinct primes dividing $(t^{a_i}-1)/(t-1)$.
The two additional parabolics are $D_{2M}$ and $Z_{2^n}:C_M$.
The identity
\[
 \gcd(M,2^{n_J}-1)=(t-1)\prod_{i\notin J}\ell_i
\]
gives all intersections. That cyclic scalar subgroup generates the full
remaining coefficient field, so its affine complement is maximal.
The four possibilities for selecting neither, either, or both soluble
types then prove (1) for every subset.

For Constructions~8 and 9 the additional group is
$K=C_{\gamma\prod_i\ell_i}:C_c$, with $c=2$ or $4$.
Compatible tori give
$K\cap S_{n_J}=C_{\gamma\prod_{i\notin J}\ell_i}:C_c$.
Lang--Steinberg gives a common representative for pairwise incident
standard subfield cosets. The Chinese remainder theorem makes
$K\longrightarrow\prod_i K/(K\cap S_{n/a_i})$ surjective as a
map of coset sets. This proves the common-representative property for
the whole family. No nonnormal coset set is treated as a quotient group.

For Construction~10, conjugate the mixed parabolic to
$Q=Z_qU_tH_t$. A prescribed pair intersection $H$ surjects onto
$Q/Z_q$, and its overgroups correspond exactly to the
$\mathbb F_t$-subspaces containing $H\cap Z_q$.
Its scalar codimension is
$\prod_i a_i-\prod_i a_i/a_j>1$.
Thus every such $H$ has a proper intermediate overgroup in $Q$.
For Construction~11, the field intersection formula and descent give
all the positive claims. If $e=1$, a final $S_1=5:4$ lies strictly
inside a torus normalizer of some $S_{a_i}$, so (1) fails.

\section{The principal correction and final remarks}
The earlier accepted principal correction uses
\[
 n=3^2\cdot5\cdot7\cdot11=3465,
 \qquad (G_1,G_2,G_3,G_4)=(S_{1155},S_{693},S_{495},S_{315}).
\]
Their common intersection is $S_3=\Sz(8)$ and every available step
is a maximal simple-subfield inclusion. All four parabolics are simple,
excluding Constructions~1--10, while the literal Construction~11 allows
rank at most three for four distinct prime divisors. This refutes
Theorem~5. The descent supplement also proves thickness, residual
connectedness and flag transitivity. The correction remains an omission
from the literal list; it is repaired by replacing that one rank bound
with $e>1$, without thereby obtaining a complete classification.

The source's introductory assertions about local two-transitivity are
quotations of the rank-three paper. That paper's complete corrected
treatment supplies the exact witnesses: two nonthin rank-three geometries
in $\Sz(8)$ and $\Sz(32)$ satisfying the literal algebraic local
condition, and a scalar-subspace family of arbitrary rank under the same
algebraic convention. The separate new coset-coordinate join now proves
that this last family consists of thick, residually connected,
flag-transitive geometries with actual local two-transitivity. Thus the
quoted rank restriction is false even for actual geometries. Its
rank-three members also occur at composite field degrees. The earlier
conservative pregeometry proof remains unchanged; this stronger
conclusion and its two-page proof are separately identified.

The diagram paper proves the conditional star statements for the surviving
prescribed families: the centre is the torus type in 8/9 and the dihedral
type in corrected 5/7; pure fields have complete diagrams. The pentagon
figure for Construction~1 presupposes an impossible lattice. The future
flag-transitivity question is answered for the canonical families 8, 9
and 11 and the retained bare-type example for 6. Flag transitivity of
corrected 5/7 and of arbitrary relative positions remains unproved here;
it was not an assertion of the source's classification theorem.

The new contribution in this assembly consists of explicit nonextension
arguments, the conditional necessity proofs for all viable rank-four
branches, the full parameter applications of the common descent and
scalar theory, the correction ledger, the diagram proof, and the
explicit join establishing actual geometries for the scalar family. Earlier
project proofs and finite certificates are retained with attribution.
The already checked Lean fragment concerns only the abstract coset-coordinate
implication; no Suzuki classification or whole manuscript is claimed
formalized. Completing this source adds no new principal-refutation count,
no ordinary recovery, no new conjecture count and no compression measurement.
\end{document}
